\documentclass[10pt,a4paper]{amsart}
\usepackage[margin=19mm]{geometry}
\usepackage{amsmath,amssymb,mathtools}
\usepackage[scr=rsfs,cal=euler]{mathalfa}
\usepackage{xfrac}
\usepackage[unicode,hidelinks,bookmarks=false]{hyperref}
\newcommand{\ie}{i.e.\ }
\newcommand{\cf}{cf.\ }
\newcommand{\resp}{respectively\ }
\newcommand{\Subsection}{\subsection}
\providecommand{\nobreakdash}{\nobreak\hbox{-}}
\setlength{\emergencystretch}{2em}
\multlinegap0pt
\usepackage{bm}
\usepackage{bbm}
\usepackage{dsfont}
\usepackage{xspace}
\usepackage{cancel}
\usepackage{mathtools}
\usepackage{amsthm}
\makeatletter
\@ifundefined{theorem}{\newtheorem{theorem}{Theorem}}{}
\@ifundefined{corollary}{\newtheorem{corollary}[theorem]{Corollary}}{}
\makeatother
\makeatletter
\newcommand{\E}{\cal{E}}
\newcommand{\FIN}{\bf{FIN}}
\newcommand{\N}{\mathbb{N}}
\renewcommand{\S}{\mathcal{S}}
\newcommand{\cal}{\mathcal}
\newcommand{\depth}{\func{depth}}
\newcommand{\dom}{\func{dom}}
\newcommand{\fin}{\mathrm{fin}}
\newcommand{\func}{\operatorname}
\newcommand{\lh}{\func{lh}}
\expandafter\ifx\csname midproof\endcsname\relax
\newenvironment{midproof}[1][Proof]
  {\begin{proof}[#1]\renewcommand{\qedsymbol}{$\blacksquare$}}
  {\end{proof}}
\fi
\newcommand{\restrictedto}{\mathord{\upharpoonright}}
\makeatother
\title[Weak A2 spaces, the Kastanas game and strategically Ramsey s]{Weak A2 spaces, the Kastanas game and strategically Ramsey sets}
\author{Clement Yung}
\date{Source: arXiv:2410.00200v4 (CC BY 4.0)}
\begin{document}
\maketitle

\noindent\emph{Source and licence.} Weak A2 spaces, the Kastanas game and strategically Ramsey sets. Version arXiv:2410.00200v4 (CC BY 4.0), distributed under the Creative Commons Attribution 4.0 International licence, as recorded in the arXiv licence metadata for this exact version and shown on the abstract page https://arxiv.org/abs/2410.00200v4. Full rights notice: licenses/arxiv-2410.00200-CC-BY-4.0.txt.\par\smallskip

\noindent\textit{Editorial scope.} Counted unit: the Corollary whose source label is \texttt{cor:V=L.gives.non.kr.set.in.examples} in the article (source0.tex lines 1123--1141, source label \texttt{cor:V=L.gives.non.kr.set.in.examples}), delivered together with the article's own complete original proof of it (source0.tex lines 1143--1180, one \texttt{proof} environment of 38 lines). No mathematics is written, repaired or re-derived: statement and proof are the article's own text line for line. Required context retained uncounted: cor:V=L.gives.non.kr.set (lines 1118--1121). Every definition, display and label of those lines is retained unchanged. Omitted from this package and not redistributed: 1--1117, 1122--1122, 1142--1142, 1181--1660. Nothing inside a retained excerpt is omitted or altered: every retained body is delivered exactly as the article's lines are written, so no omission receipt of any kind accompanies this package. Citations: the retained excerpts carry no citation command at all, so no bibliography entry of the article is transcribed and no reference is redistributed. Reference adaptation: none was needed and none was made; every reference inside the delivered unit resolves inside it, so no unresolved reference or citation is produced. Printed slips kept literal: no printed word, symbol, hypothesis or number of the counted statement or of its proof was altered. Layout and class: the article's own class is \texttt{amsart} with packages including amsmath, amsthm, amssymb, amsfonts, bbm, bm, graphicx, xcolor, hyperref, enumitem, cleveref, biblatex; the wrapper is the lane's pinned \texttt{amsart} editorial wrapper at 10pt on A4, which restates the theorem-like environments the retained text needs and adds the article's own macro definitions that the retained text needs (\texttt{\textbackslash E}, \texttt{\textbackslash FIN}, \texttt{\textbackslash N}, \texttt{\textbackslash S}, \texttt{\textbackslash cal}, \texttt{\textbackslash depth}, \texttt{\textbackslash dom}, \texttt{\textbackslash fin}, \texttt{\textbackslash func}, \texttt{\textbackslash lh}, \texttt{\textbackslash restrictedto}), with extra wrapper packages (bm, bbm, dsfont, xspace, cancel, mathtools). The delivered numbering is the unit's own. Assets: no retained span contains a figure environment, a graphics-inclusion command, a PDF-page inclusion, a table or a TikZ picture; no image file is copied into this package, the package declares no asset dependency and no image file is redistributed with it. Licence: version arXiv:2410.00200v4 of this article is distributed under the Creative Commons Attribution 4.0 International licence, as recorded in the arXiv licence metadata for this exact version and shown on the abstract page https://arxiv.org/abs/2410.00200v4. A full rights notice is delivered at licenses/arxiv-2410.00200-CC-BY-4.0.txt. Characters: every retained excerpt is pure ASCII, so the delivered engine, XeLaTeX with its default Latin Modern text font, reports no missing character.
\begin{corollary}[$\mathsf{V=L}$]
\label{cor:V=L.gives.non.kr.set}
    Let $(\cal{R},\leq,r)$ be a deep \textbf{wA2}-space, and assume that $\cal{AR}$ is countable. Then there exists a $\mathbf{\Sigma}_2^1$ subset of $\cal{R}$ which is not Kastanas Ramsey.
\end{corollary}

\begin{corollary}[$\mathsf{V=L}$]
\label{cor:V=L.gives.non.kr.set.in.examples}
    The following \textbf{wA2}-spaces have a $\mathbf{\Sigma}_2^1$ subset which is not Kastanas Ramsey:
    \begin{enumerate}
        \item $([\N]^\infty,\subseteq,r)$.

        \item $(\FIN_k^{[\infty]},\leq,r)$, the topological Ramsey space of infinite block sequences.
        
        \item $(\FIN_{\pm k}^{[\infty]},\leq,r)$, a variant of the space of infinite block sequences.

        \item $(W_{Lv}^{[\infty]},\leq,r)$, the Hales-Jewett space.

        \item $(\mathcal{S}^\infty,\subseteq,r)$, the topological Ramsey space of strong subtrees.

        \item $(\E^\infty,\leq,r)$, the Carlson-Simpson space.

        \item $(E^{[\infty]},\leq,r)$, the space of infinite-dimensional block subspaces of a countable vector space.
    \end{enumerate}
\end{corollary}

\begin{proof}
    By Corollary \ref{cor:V=L.gives.non.kr.set}, it suffices to show that every \textbf{wA2}-space above is deep.
    \begin{enumerate}
        \item Let $A = \{n_0,n_1,\dots\} \in [\N]^\infty$, and let $a \subseteq A$ be finite. For all $N$ such that $n_N > \max(a)$, we have that $\depth_A(a \cup \{n_{N-1}\}) = N$. Therefore, $a \cup \{n_{N-1}\} \in r_{\lh(a)+1}[a,A]$ and $\depth_A(a \cup \{n_{N-1}\}) \geq N$.
        
        \item Let $A = (x_0,x_1,\dots) \in \FIN_k^{[\infty]}$, and let $a \in \FIN_k^{[<\infty]}\restrictedto A$. For all $N$ such that $x_N > a$, we have that $\depth_A(a^\frown x_{N-1}) = N$. Therefore, $a^\frown x_{N-1} \in r_{\lh(a)+1}[a,A]$ and $\depth_A(a^\frown x_{N-1}) \geq N$.
        
        \item The proof is identical to that of $(\FIN_k^{[\infty]},\leq,r)$.
        
        \item Let $A = (x_0,x_1,\dots) \in W_{Lv}^{[\infty]}$, and let $a = (y_i)_{i<n} \in W_{Lv}^{[<\infty]}\restrictedto A$. Since $A$ is rapidly increasing, for $N$ large enough we have that $\sum_{i<n} |y_i| < |x_N|$. Therefore, $a^\frown x_N \in r_{\lh(a)+1}[a,A]$ and $\depth_A(a^\frown x_N) \geq N$.
        
        \item Let $A \subseteq T$ be a strong subtree. Given any $a \in \S_{<\infty}\restrictedto A$, we let $S_a$ be the set of terminal nodes in $a$. Since $a$ is a strong subtree, $S_a \subseteq \func{split}_N(A)$ for some $N$. We observe that:
        \begin{align*}
            \depth_A(a) = N \iff S_a \subseteq \func{split}_N(A).
        \end{align*}
        Fix any $a \in \S_{<\infty}\restrictedto A$ and $N < \omega$ be large enough. For each $s \in S_a$, let $t_{s,0},t_{s,1} \in \func{split}_N(A)$ be such that $s \sqsubseteq t_{s,0}$, $s \sqsubseteq t_{s,1}$ and $t_{s,0} \neq t_{s,1}$. We let:
        \begin{align*}
            b := \{u \in A : u \sqsubseteq t_{s,i} \text{ for some $s \in S_a$ and $i \in \{0,1\}$}\}.
        \end{align*}
        Observe that $a$ is an initial segment of $b$, and every terminal node in $a$ splits in $b$. Thus, $b \in r_{\lh(a)+1}[a,A]$ and $\depth_A(b) = N$.

        \item Let $A \in \E_\infty$, and let $a \in \cal{AE}_\infty\restrictedto A$. Let $\{p_n(A) : n < \omega\}$ be the increasing enumeration of the minimal representatives of $A$. Given $N > m := \depth_A(a)$, we define the equivalence relation $b$ on $\{0,1,\dots,p_N(A)-1\}$ as follows: Given $i,j \in \N$, we define:
        \begin{small}
            \begin{gather*}
                (i,j) \in b \iff
                \begin{cases}
                    (i,j) \in A \text{ and } (i \in \dom(a) \text{ or } j \in \dom(a)), \text{ or}; \\
                    (i \notin \dom(a) \text{ and } j \notin \dom(a)).
                \end{cases}
            \end{gather*}
        \end{small}
        We shall show that $b \in r_{\lh(a)+1}[a,A]$ and $\depth_A(b) = N$.
        
        Given $i,j < p_N(A)$ such that $(i,j) \in A$, if $i \in \dom(a)$ or $j \in \dom(a)$ then $(i,j) \in b$. Otherwise, $(i,j) \in b$ as well. Therefore, $b$ is an equivalence relation on $\dom(r_N(A))$ which is coarser than $A$, so $b \leq_\fin r_N(A)$. To see that $a \sqsubseteq b$ - if $i,j \in \dom(a)$ and $(i,j) \in b$, then $(i,j) \in A$, so $(i,j) \in a$ as $a$ is coarser than $A$. Finally, the equivalence classes in $b$ are either of the form $[i]_b$ for some $i \in \dom(a)$ (of which $(i,j) \notin a$ implies that $[i]_b \neq [j]_b$), or $[i]_b$ for any $i \notin \dom(a)$ (of which $[i]_b = [j]_b$ for all $i,j \notin \dom(a)$). Therefore, $b$ has $\lh(a) + 1$ many equivalence classes, i.e. $b \in r_{\lh(a)+1}[a,A]$.

        \item The proof is identical to that of $(\FIN_k^{[\infty]},\leq,r)$.
    \end{enumerate}
\end{proof}

\end{document}
